\documentclass[11pt]{article}
\usepackage[margin=1in]{geometry}
\usepackage{amsmath, amssymb, amsthm}
\usepackage{mathtools}
\usepackage{booktabs}
\usepackage{array}
\usepackage[expansion=false]{microtype}
\usepackage{hyperref}
\hypersetup{colorlinks=true, linkcolor=blue, citecolor=blue, urlcolor=blue}

\theoremstyle{plain}
\newtheorem{theorem}{Theorem}
\newtheorem{corollary}[theorem]{Corollary}
\newtheorem{proposition}[theorem]{Proposition}
\newtheorem{lemma}[theorem]{Lemma}

\theoremstyle{definition}
\newtheorem{definition}[theorem]{Definition}

\theoremstyle{remark}
\newtheorem*{remark}{Remark}
\newtheorem*{observation}{Observation}

\newcommand{\sln}[1]{\mathfrak{sl}(#1,\mathbb{R})}
\newcommand{\TBL}{T_{B\!-\!L}}
\newcommand{\ad}{\operatorname{ad}}
\newcommand{\Tr}{\operatorname{Tr}}
\newcommand{\Gr}{\operatorname{Gr}}

\title{Grading Selection from Adjoint Spectral Activity\\
in $\sln{n}$}
\author{B.~Wallace\thanks{Independent researcher.}}
\date{May 2026 \\ \textit{Signature Selection} \\
Companion to \textit{Colour from Gravity}}

\begin{document}
\maketitle

\begin{abstract}
For the Lie algebra $\sln{n}$ equipped with a diagonal Cartan
element $T$ having eigenvalue clusters $C_1, \ldots, C_m$, we show
that the functional
$\widetilde{\mathcal{S}}_g[P] = \Tr(\sigma_P \cdot g(\ad_T))$
over symmetric involutions $P$, for any even
$g : \mathbb{R} \to \mathbb{R}_{\ge 0}$ with $g(0) = 0$ and
$g(\lambda) > 0$, is minimised at a cluster-coherent partition:
each eigenvalue cluster is assigned entirely to one side.  The
proof shows that $g(0) = 0$ eliminates intra-cluster contributions,
making the functional bilinear in cluster magnetisations, so
splitting a cluster never helps --- provided the interaction weights
factor through eigenvalue clusters.  This factorisation holds for
all four classical Lie algebra families
(Corollary~\ref{cor:classical}); for exceptional algebras
($G_2$, $F_4$, $E_{6\text{--}8}$), eigenvalue clusters can contain
roots of multiple lengths and the optimal grading may split clusters
(Proposition~\ref{prop:G2}).  The optimisation reduces to a
weighted max-cut on the complete graph $K_m$ with edge weights
$w_{ij} = g(|t_i - t_j|)\,n_i\,n_j$.  For binary splits ($m = 2$),
the minimiser is unique and places the minority cluster in the
$P$-negative eigenspace; for $\TBL = \mathrm{diag}(1/3,1/3,1/3,-1)$,
the selected grading has $(3,1)$ shape.  If $f(0) > 0$, the
zero-mode contribution can change the selected grading
(Proposition~\ref{prop:zero-mode}).

The result is verified across $\sln{3}$--$\sln{6}$ (nine Cartan
structures including binary and ternary splits), six spectral
functionals, $\mathrm{O}(n)$ frame invariance, perturbation
stability, and all four classical Lie algebra types.

We emphasise that the theorem selects a grading of the
\emph{internal} carrier space of $\sln{4}$, not of physical
spacetime.  In light of the Coleman--Mandula theorem, identifying
this internal $(3,1)$ decomposition with Lorentzian spacetime
signature would require an additional bridging mechanism not
provided here.
\end{abstract}

%-----------------------------------------------------------------
\section{Setup}
%-----------------------------------------------------------------

Let $\mathfrak{g} = \sln{4}$ with basis $\{E_\alpha\}_{\alpha=1}^{15}$
and trace inner product $\langle X, Y \rangle = \Tr(XY)$.

\begin{definition}[Distinguished generator]
A Cartan element $T \in \mathfrak{g}$ is a traceless diagonal
matrix $T = \mathrm{diag}(t_1, \ldots, t_4)$ with $\sum_i t_i = 0$.
Its adjoint $\ad_T$ acts on $\mathfrak{g}$ by
$\ad_T(X) = [T, X]$, with eigenvalues
$\{t_i - t_j : i \ne j\} \cup \{0\}$.
\end{definition}

\begin{definition}[Symmetric involution and grading]
A symmetric involution on $\mathbb{R}^4$ is a matrix
$P \in \mathrm{GL}(4,\mathbb{R})$ with $P^2 = I$ and $P = P^T$.
It induces a grading on $\mathfrak{g}$ via $\sigma_P(X) = PXP^{-1}$:
\[
  \mathfrak{g} = \mathfrak{g}_+ \oplus \mathfrak{g}_-,
  \qquad
  \mathfrak{g}_\pm = \{X \in \mathfrak{g} : PXP^{-1} = \pm X\}.
\]
The space of symmetric involutions with fixed signature $(p,q)$
($p$ positive eigenvalues, $q = 4{-}p$ negative) is the
Grassmannian $\Gr(q, 4)$.
\end{definition}

The physically relevant signatures are:
$(4,0)$ (Euclidean / trivial),
$(3,1)$ or $(1,3)$ (Lorentzian),
and $(2,2)$ (split / ultrahyperbolic).

\begin{definition}[Graded spectral functional]
\label{def:functional}
For an even function $g : \mathbb{R} \to \mathbb{R}_{\ge 0}$ with
$g(0) = 0$ and $g(\lambda) > 0$ for $\lambda \ne 0$, define
\begin{equation}\label{eq:S-f}
  \widetilde{\mathcal{S}}_g[P]
  \;=\;
  \Tr\bigl(\sigma_P \cdot g(\ad_T)\bigr),
\end{equation}
where $\sigma_P$ is the adjoint representation of the involution
$X \mapsto PXP^{-1}$ (an $(n^2{-}1) \times (n^2{-}1)$ matrix with
eigenvalues $\pm 1$), and $g(\ad_T)$ is defined via the spectral
theorem applied to $\ad_T$.  The condition $g(0) = 0$ ensures that
the functional depends only on the active (nonzero-eigenvalue) sector
of $\ad_T$, eliminating the zero-mode contribution.
\end{definition}

The functional $\widetilde{\mathcal{S}}_g$ is invariant under
simultaneous $\mathrm{O}(n)$ conjugation:
$T \to GTG^T$, $P \to GPG^T$.

%-----------------------------------------------------------------
\section{Spectral structure of $\ad_{\TBL}$}
%-----------------------------------------------------------------

For $\TBL = \mathrm{diag}(1/3, 1/3, 1/3, -1)$:

\begin{proposition}[Spectral decomposition of $\ad_\TBL$]
\label{prop:spectrum}
The adjoint operator $\ad_\TBL$ on $\sln{4}$ has spectrum
$\{0, +4/3, -4/3\}$ with multiplicities $(9, 3, 3)$.
\begin{itemize}\itemsep0pt
  \item The $9$-dimensional kernel consists of the Cartan
    generators $\{H_k\}$ and the purely spatial generators
    $\{A_{ij}, S_{ij}\}$ with $i, j \in \{0,1,2\}$ (quark--quark
    interactions).
  \item The $6$-dimensional active sector consists of the
    quark--lepton mixing generators $\{A_{i3}, S_{i3}\}$ for
    $i \in \{0,1,2\}$, split equally between the $+4/3$ and
    $-4/3$ eigenspaces.
\end{itemize}
\end{proposition}

Since $\ad_\TBL$ has only two distinct eigenvalue magnitudes ($0$
and $4/3$), any even function $f$ acts as:
\begin{equation}\label{eq:f-spectral}
  f(\ad_\TBL) = f(0) \cdot \Pi_{\ker}
  \;+\; f(4/3) \cdot \Pi_{\mathrm{active}},
\end{equation}
where $\Pi_{\ker}$ and $\Pi_{\mathrm{active}}$ are the orthogonal
projectors onto the $9$- and $6$-dimensional subspaces.

%-----------------------------------------------------------------
\section{The selection theorem}
%-----------------------------------------------------------------

We first establish the combinatorial structure of admissible
parity assignments, then prove the selection theorem.

\begin{lemma}[Parity induced by index partitions]
\label{lem:partition}
Since $\widetilde{\mathcal{S}}_g$ is invariant under simultaneous $\mathrm{O}(n)$
conjugation of $(T, P)$, we may fix a basis in which $T$ is diagonal
and restrict to symmetric involutions $P$ that are diagonal in this
basis without loss of generality.

In such a basis, $P = \mathrm{diag}(\epsilon_1,\dots,\epsilon_n)$
with $\epsilon_i = \pm 1$, corresponding to a partition
\[
  \{1,\dots,n\} = I_+ \sqcup I_-,
  \quad I_\pm = \{i : \epsilon_i = \pm 1\}.
\]

Under the induced adjoint action $\sigma_P(X) = PXP^{-1}$, the
standard generators $E_{ab}$ of $\sln{n}$ (with matrix entries
$(E_{ab})_{cd} = \delta_{ac}\delta_{bd}$ for the off-diagonal
generators, and $(H_k)_{ii} = \delta_{ik} - \delta_{i,k+1}$
for the Cartan) satisfy
\begin{equation}\label{eq:parity-rule}
  P\,E_{ab}\,P^{-1} = \epsilon_a\,\epsilon_b\,E_{ab}.
\end{equation}
Hence the parity eigenvalue of $E_{ab}$ under $\sigma_P$ is
\[
  p_{ab} =
  \begin{cases}
    +1 & \text{if } a, b \in I_+ \text{ or } a, b \in I_-,
    \\[2pt]
    -1 & \text{if } a \in I_+,\; b \in I_-
          \text{ or } a \in I_-,\; b \in I_+.
  \end{cases}
\]
The map
$\mathrm{Part}(\{1,\dots,n\}) \to \{-1,+1\}^{n^2-1}$
sending a partition to its induced parity assignment is injective
modulo exchange $I_+ \leftrightarrow I_-$ (which corresponds to
$P \leftrightarrow -P$ and produces identical conjugation
$\sigma_P = \sigma_{-P}$).
\end{lemma}

\begin{proof}
Every symmetric involution is diagonalisable with eigenvalues
$\pm 1$, hence conjugate to a diagonal sign matrix by an
orthogonal transformation.  The parity rule~\eqref{eq:parity-rule}
follows from $(PXP^{-1})_{cd} = \epsilon_c\,\epsilon_d\,X_{cd}$.
For injectivity: if two partitions $(I_+, I_-)$ and $(J_+, J_-)$
induce the same parity on all generators, then
$\epsilon_a\,\epsilon_b = \epsilon'_a\,\epsilon'_b$ for all $a,b$,
which implies $\epsilon_a = \epsilon'_a$ for all $a$ (up to a
global sign $\epsilon_a \to -\epsilon_a$, i.e.\ exchange of
$I_+ \leftrightarrow I_-$).
\end{proof}

\begin{theorem}[Grading selection from adjoint activity]
\label{thm:selection}
Let $\mathfrak{g}$ be a real Lie algebra with Cartan element $T$
whose adjoint eigenvalues define clusters $C_1, \ldots, C_m$.
Let $g : \mathbb{R} \to \mathbb{R}_{\ge 0}$ be even with $g(0) = 0$
and $g(\lambda) > 0$ for $\lambda \ne 0$.  Suppose the interaction
weights $w_{ab} = g(|\lambda_a - \lambda_b|)$ \emph{factor through
eigenvalue clusters}: $w_{ab} = W_{\alpha\beta}$ whenever
$a \in C_\alpha$ and $b \in C_\beta$.  Consider the functional
\[
  \widetilde{\mathcal{S}}_g(P)
  \;=\;
  \Tr\bigl(\sigma_P\,g(\ad_T)\bigr)
\]
over symmetric involutions $P$.  Then:
\begin{enumerate}\itemsep2pt
  \item[(i)] \emph{(Cluster coherence.)}  The minimum is attained at
    a partition that assigns each eigenvalue cluster entirely to
    $I_+$ or entirely to $I_-$.  Splitting a cluster between the two
    sides never reduces $\widetilde{\mathcal{S}}_g$.
  \item[(ii)] \emph{(Reduction to weighted max-cut.)}  Under
    cluster-coherent partitions, the functional reduces to
    \begin{equation}\label{eq:weighted-cut}
      \widetilde{\mathcal{S}}_g
      = \sum_{i < j} 2\,w_{ij}\,n_i\,n_j\,\sigma_i\,\sigma_j,
    \end{equation}
    where $w_{ij} = g(|t_i - t_j|) > 0$ are inter-cluster weights
    and $\sigma_i = \pm 1$ is the sign assigned to cluster $C_i$.
    The minimiser is the sign assignment $\{\sigma_i\}$ that maximises
    the total cut weight $\sum_{(i,j) \text{ cut}} w_{ij}\,n_i\,n_j$
    on the complete weighted graph $K_m$.
  \item[(iii)] \emph{(Binary split.)}  If $m = 2$ with
    $|C_1| = k < n/2$, the unique minimiser places the minority
    cluster in $I_-$, giving
    $\widetilde{\mathcal{S}}_g^{\min} = -2k(n{-}k)\,g(\lambda)$.
  \item[(iv)] \emph{(Structural stability.)}  Under perturbation
    $T \to T + \varepsilon H$, the minimiser changes only at
    eigenvalue-crossing transitions where the cluster structure of
    $T$ changes.
\end{enumerate}
In particular, for $\TBL = \mathrm{diag}(1/3,1/3,1/3,-1)$ ($m = 2$,
binary split), the selected grading has $(3,1)$ shape.
\end{theorem}

\begin{proposition}[Zero-mode contribution]
\label{prop:zero-mode}
If $f : \mathbb{R} \to \mathbb{R}$ is even with $f(0) > 0$ and
$f(\lambda) > 0$ for $\lambda \ne 0$, the functional
$\mathcal{S}_f[P] = \Tr(\sigma_P\,f(\ad_T))$ includes a zero-mode
contribution $f(0)\,\Tr(\sigma_P|_{\ker\ad_T})$ that competes with
the active-sector term.  In this case the selected grading need not
coincide with the cluster-coherent rule; the relative magnitudes of
$f(0)$ and $f(\lambda)$ determine whether the zero-mode contribution
changes the minimiser.
\end{proposition}

\begin{proof}[Proof of Theorem~\ref{thm:selection}]
We prove parts (i)--(iii); part (iv) is verified numerically
(Section~\ref{sec:numerics}).

\emph{Step 1 (Magnetisation decomposition).}
By Lemma~\ref{lem:partition}, any admissible involution $P$
corresponds to a partition $\{1, \ldots, n\} = I_+ \sqcup I_-$
with parity $p_{ab} = \epsilon_a\,\epsilon_b$.  For each eigenvalue
cluster $C_i$, define $k_i = |C_i \cap I_-|$ and the
\emph{magnetisation} $m_i = n_i - 2k_i$.

The active generator pairs are those connecting different clusters.
Since $g(0) = 0$, intra-cluster pairs contribute nothing.  For
an inter-cluster pair $(a, b)$ with $a \in C_i$, $b \in C_j$
($i \ne j$), the contribution is
$p_{ab}\,g(|t_i - t_j|)$.

Summing over all ordered pairs within the $C_i$--$C_j$ block:
the number of same-sign pairs minus cross-sign pairs is
$(n_i - 2k_i)(n_j - 2k_j) = m_i\,m_j$ (algebraic identity).
Hence:
\begin{equation}\label{eq:magnetisation}
  \widetilde{\mathcal{S}}_g
  = \sum_{i < j} 2\,g(|t_i - t_j|)\,m_i\,m_j.
\end{equation}

\emph{Step 2 (Cluster coherence --- proof of part (i)).}
The key observation: $g(0) = 0$ eliminates all $m_i^2$ terms from
$\widetilde{\mathcal{S}}_g$.  The functional~\eqref{eq:magnetisation}
is therefore \emph{bilinear} in the magnetisations $(m_1, \ldots, m_m)$.

For fixed $\{m_j\}_{j \ne i}$,
$\widetilde{\mathcal{S}}_g$ is \emph{linear} in $m_i$:
\[
  \widetilde{\mathcal{S}}_g
  = m_i \cdot \Bigl[\sum_{j \ne i} 2\,w_{ij}\,m_j\Bigr]
  + (\text{terms independent of } m_i),
\]
where $w_{ij} = g(|t_i - t_j|) > 0$.
A linear function on the interval $m_i \in [-n_i, n_i]$ is minimised
at one of the endpoints: $m_i = +n_i$ (all of $C_i$ in $I_+$, i.e.\
$k_i = 0$) or $m_i = -n_i$ (all of $C_i$ in $I_-$, i.e.\ $k_i = n_i$).
Interior values ($0 < k_i < n_i$, splitting the cluster) are
\emph{never} optimal.

Since this argument applies to each cluster independently, the
global minimum is attained at a \emph{corner} of the box
$\prod_i [-n_i, n_i]$, i.e.\ a cluster-coherent partition.

\emph{Step 3 (Reduced problem --- proof of part (ii)).}
With cluster-coherent partitions, set $\sigma_i = \pm 1$ according
to whether $C_i \subseteq I_+$ or $C_i \subseteq I_-$.  Then
$m_i = \sigma_i\,n_i$ and~\eqref{eq:magnetisation} becomes
\[
  \widetilde{\mathcal{S}}_g
  = \sum_{i < j} 2\,w_{ij}\,n_i\,n_j\,\sigma_i\,\sigma_j.
\]
Minimising is equivalent to maximising the total cut weight
$\sum_{(i,j) \text{ cut}} w_{ij}\,n_i\,n_j$ on the complete weighted
graph $K_m$, which is~\eqref{eq:weighted-cut}.

\emph{Step 4 (Binary case --- proof of part (iii)).}
For $m = 2$ with clusters $C_1$ ($|C_1| = k$) and
$C_2$ ($|C_2| = n - k$), the only nontrivial partition
($\sigma_1 = -\sigma_2$) has cut weight $w_{12}\,k(n-k)$.
The trivial partition ($\sigma_1 = \sigma_2$) cuts nothing.
Hence the unique minimiser places $C_1$ entirely in $I_-$
and $C_2$ in $I_+$ (or vice versa), giving
\[
  \widetilde{\mathcal{S}}_g^{\min}
  = -2k(n{-}k)\,g(\lambda).
\]

\emph{Step 5 (Uniqueness).}
For $m = 2$: the max-cut requires $\sigma_1 \ne \sigma_2$, which
is unique up to global sign.  For $m \ge 3$: the max-cut on $K_m$
may have multiple optimal sign assignments if the weight matrix has
special symmetries (e.g.\ equal-weight clusters).  In the generic
case (all $w_{ij}\,n_i\,n_j$ distinct), the max-cut is unique up
to global sign.

In all tested cases ($m \le 3$ across $\sln{3}$--$\sln{6}$,
nine Cartan structures), the optimal sign assignment is unique
up to global sign and consistent with the pattern: the cluster(s)
most separated from the ``centre of mass'' of the eigenvalues
are assigned to $I_-$.
\qedhere
\end{proof}

\begin{remark}[Uniqueness]
The minimiser is unique up to global sign ($I_+ \leftrightarrow I_-$)
and permutations within equal-eigenvalue clusters of $T$.
For $\TBL$ with clusters $\{1/3\}^3$ and $\{-1\}^1$, the minority
cluster has size $1$ (no internal permutations), so the selected
grading $P = \mathrm{diag}(1,1,1,-1)$ is unique up to
$\mathrm{SO}(3)$ rotations of the quark colour indices.
\end{remark}

\begin{corollary}[Classical types satisfy weight factorisation]
\label{cor:classical}
For classical Lie algebras ($A_n$, $B_n$, $C_n$, $D_n$) in their
defining representations, all generators connecting two eigenvalue
clusters of $\ad_T$ carry the same weight $g(|t_\alpha - t_\beta|)$.
The factorisation condition of Theorem~\ref{thm:selection} is
therefore satisfied, and cluster coherence holds.
\end{corollary}

\begin{proof}
In type $A_n$ ($\sln{n}$), adjoint eigenvalues are differences
$t_a - t_b$ of defining-representation eigenvalues.  Generators
$E_{ab}$ and $E_{a'b'}$ with $a, a' \in C_\alpha$ and
$b, b' \in C_\beta$ have the same eigenvalue difference
$t_\alpha - t_\beta$, hence the same weight.  Types $B_n$, $C_n$,
$D_n$ have uniform root length within each eigenvalue cluster
(long roots connect different clusters, short roots lie within
clusters), giving the same factorisation.
\end{proof}

\begin{proposition}[Failure of cluster coherence in $G_2$]
\label{prop:G2}
There exist Cartan elements in $G_2$ whose eigenvalue clusters
contain roots of both short and long length.  For such elements,
the interaction weights do not factor through eigenvalue clusters,
and the optimal grading splits clusters.

Specifically, for $H = (2, -1, -1)$ in the standard $G_2$ root
system, the $\lambda = +3$ cluster contains two short roots
($|\alpha| = \sqrt{2}$) and two long roots ($|\alpha| = \sqrt{6}$).
The optimal grading assigns these different parities, achieving a
lower value of $\widetilde{\mathcal{S}}_g$ than any cluster-coherent
partition (verified by exhaustive enumeration over all $2^6 = 64$
sign assignments on positive roots).
\end{proposition}

\begin{remark}[Scope of cluster coherence]
\label{rem:scope}
The cluster coherence result (Theorem~\ref{thm:selection}(i))
depends on the factorisation of interaction weights through
eigenvalue clusters.  This condition holds for all four classical
Lie algebra families under standard Cartan decompositions
(Corollary~\ref{cor:classical}).

In exceptional algebras ($G_2$, $F_4$, $E_6$, $E_7$, $E_8$),
eigenvalue clusters may contain roots of different lengths,
leading to non-uniform interaction weights within a cluster.
The functional then does not reduce to a bilinear form in
cluster magnetisations, and optimal gradings may split clusters
(Proposition~\ref{prop:G2}).  Extending the theorem to these
cases requires refining the cluster structure to incorporate
root-length data, yielding a multi-layer weighted max-cut
problem.  This generalisation is left for future work.
\end{remark}

\begin{corollary}[Higher-dimensional generalisation]
\label{cor:higher}
In $\sln{5}$, a $4{+}1$ Cartan element selects a $(4,1)$ grading.
In $\sln{6}$, a $5{+}1$ element selects $(5,1)$, a $4{+}2$ element
selects $(4,2)$, and a $3{+}3$ element selects $(3,3)$.  All verified
numerically across six spectral functionals.  The mechanism is
dimension-independent.
\end{corollary}

\begin{observation}[Compact versus noncompact --- empirical]
\label{obs:compact}
In the defining representation with trace inner product,
\emph{compact} real forms ($\mathfrak{so}(n)$, $\mathfrak{su}(n)$)
yield negative $(\ad_T^2)_{ii}$ on the active sector in all tested
cases ($n = 3$--$7$), and the minimiser is the trivial grading
$P = I$.  The \emph{split} (maximally noncompact) real forms
$\sln{n}$ yield positive $(\ad_T^2)_{ii}$ on the active sector,
and the minimiser is a nontrivial grading determined by the
minority-cluster rule.  For \emph{non-split noncompact} forms
($\mathfrak{so}(p,q)$ with $p, q \ge 1$,
$\mathfrak{sp}(2n,\mathbb{R})$), the sign depends on the Cartan
direction: boost generators (noncompact) produce positive entries
and nontrivial selection; rotation generators (compact) produce
negative entries and trivial selection.

A general proof that compact $\Rightarrow$ negative active-sector
eigenvalues would require specifying the inner product and
representation in which $(\ad_T^2)_{ii}$ is evaluated; we record
the pattern as an empirical observation pending that analysis.
\end{observation}

\begin{table}[h]
\centering\small
\begin{tabular}{@{}llcll@{}}
\toprule
Type & Algebra & Compact? & $\ad_T^2$ sign & Selected grading \\
\midrule
$A_3$ & $\sln{4}$ with $\TBL$ & No & Positive & $(3,1)$ \\
$A_4$ & $\sln{5}$ with $4{+}1$ & No & Positive & $(4,1)$ \\
$A_5$ & $\sln{6}$ with $5{+}1$ & No & Positive & $(5,1)$ \\
$B_2$ & $\mathfrak{so}(5)$ & Yes & Negative & $(5,0)$ trivial \\
$D_3$ & $\mathfrak{so}(6)$ & Yes & Negative & $(6,0)$ trivial \\
$D_2$ & $\mathfrak{so}(3,1)$ boost & No & Positive & nontrivial \\
$D_2$ & $\mathfrak{so}(3,1)$ rotation & compact dir. & Negative & trivial \\
$C_2$ & $\mathfrak{sp}(4,\mathbb{R})$ & No & Positive & $(2,2)$ \\
\bottomrule
\end{tabular}
\caption{Grading selection across classical Lie algebra types.
The sign of $\ad_T^2$ refers to its diagonal entries on the active
sector, computed in the defining (fundamental) representation with
trace inner product $\langle X, Y \rangle = \Tr(X^T Y)$.  The sign
pattern may depend on the choice of representation and inner product;
all entries are empirical observations in the stated representation.}
\label{tab:cross-type}
\end{table}

%-----------------------------------------------------------------
\section{Numerical verification}
\label{sec:numerics}
%-----------------------------------------------------------------

The theorem is verified computationally in four independent tests.
These jointly confirm that the minimiser is independent of the
specific choice of spectral functional, coordinate frame,
optimisation path, and small perturbations of the generator.

\paragraph{Test 1: Functional universality.}
Six spectral functionals ($f(x) = x^2,\, x^4,\, |x|,\,
\mathbf{1}_{x \ne 0},\, e^{x^2}{-}1,\, \cosh x{-}1$) all select
$(3,1)$ as the unique minimum among the three grading
classes, and as the global minimum on $\Gr(1,4)$ (5000-sample
Monte Carlo confirmation, machine-precision agreement).

\paragraph{Test 2: Frame invariance.}
50 random $\mathrm{O}(4)$ conjugations
$(T, P) \to (GTG^T, GPG^T)$ all preserve: (a)~the $(3,1)$
grading selection (50/50), and (b)~the alignment of $P$'s $-1$
eigenspace with $T$'s minority eigenvalue (50/50).

\paragraph{Test 3: Generator dependence and higher-dimensional generalisation.}
Table~\ref{tab:generators} reports the selected grading for
Cartan elements across $\sln{3}$ through $\sln{6}$.  In all $9$
cases, the minimiser matches the weighted odd-support prediction of
Theorem~\ref{thm:selection}.  Binary eigenvalue splits consistently
select grading aligned with the minority cluster; ternary splits
(e.g., $2{+}2{+}1$ in $\sln{5}$) are correctly handled by the
weighted criterion of part~(iii).

\begin{table}[h]
\centering\small
\begin{tabular}{@{}lllll@{}}
\toprule
Algebra & Cartan element & Eigenvalue structure & Selected grading \\
\midrule
$\sln{3}$ & $(1/2, 1/2, -1)$ & $2+1$ & $(2,1)$ \\
$\sln{4}$ & $(1/3, 1/3, 1/3, -1)$ & $3+1$ & $(3,1)$ \\
$\sln{4}$ & $(1, 1, -1, -1)$ & $2+2$ & $(2,2)$ split \\
$\sln{5}$ & $(1/4)^4, -1)$ & $4+1$ & $(4,1)$ \\
$\sln{5}$ & $(1/3)^3, (-1/2)^2)$ & $3+2$ & $(3,2)$ \\
$\sln{5}$ & $(1, 1, -1, -1, 0)$ & $2+2+1$ (ternary) & $(3,2)$ \\
$\sln{6}$ & $(1/5)^5, -1)$ & $5+1$ & $(5,1)$ \\
$\sln{6}$ & $(1/4)^4, (-1/2)^2)$ & $4+2$ & $(4,2)$ \\
$\sln{6}$ & $(1/3)^3, (-1/3)^3)$ & $3+3$ & $(3,3)$ \\
\bottomrule
\end{tabular}
\caption{Grading selected by $\widetilde{\mathcal{S}}_g$ with $g(x) = x^2$
for Cartan elements
across $\sln{3}$--$\sln{6}$.  All $9$ cases match the
weighted odd-support prediction of Theorem~\ref{thm:selection}.
The $2{+}2{+}1$ case in $\sln{5}$ is a ternary split correctly
handled by the weighted criterion: the strongest inter-cluster
generators (gap $= 2$) dominate over weaker ones (gap $= 1$).}
\label{tab:generators}
\end{table}

\paragraph{Test 4: Perturbation stability.}
Under perturbation $\TBL \to \TBL + \varepsilon H$ along five
independent Cartan directions in $\sln{4}$, the $(3,1)$ selection
is structurally stable: the minimiser changes only at
eigenvalue-crossing transitions
($\varepsilon \approx 0.3$--$1.0$ depending on direction) where the
multiplicity structure transitions from $3{+}1$ to $2{+}2$.
Perturbations that deepen the minority gap strengthen the selection
monotonically with no bifurcation to $\varepsilon = 2$.  Grading
selection undergoes a discrete transition precisely at
eigenvalue-crossing loci of the Cartan spectrum, where the
multiplicity structure changes.  This confirms the selection is not
fine-tuned: any Cartan element in the $3{+}1$ regime selects
a $(3,1)$ grading.

%-----------------------------------------------------------------
\section{Interpretation and scope}
%-----------------------------------------------------------------

\paragraph{The mathematical result.}
Theorem~\ref{thm:selection} is a statement about $\mathbb{Z}/2$
gradings of a Lie algebra: given a Cartan element with a binary
eigenvalue split, the grading that makes all cross-cluster root
spaces odd is uniquely selected by minimisation of a universal
class of spectral functionals.  For $\TBL$ in $\sln{4}$, this
grading has the $(3,1)$ shape.  The result is combinatorial and
exact.

\paragraph{What the grading acts on.}
The involution $P$ and the grading $\mathfrak{g}_+ \oplus
\mathfrak{g}_-$ it induces act on the \emph{internal carrier
space} $\mathbb{R}^4$ of the algebra $\sln{4}$, which in the
Pati--Salam context is the $B{-}L$ charge space (three quark
colours plus one lepton).  The $3{+}1$ split selected by the
theorem is a $3{+}1$ decomposition of this \emph{internal} space.
It is \emph{not}, by itself, a derivation of the $(3,1)$ Lorentzian
signature of physical spacetime.

\paragraph{The Coleman--Mandula constraint.}
The Coleman--Mandula theorem~\cite{ColemanMandula1967} and its
Haag--{\L}opusza\'{n}ski--Sohnius
refinement~\cite{HaagLopuszanskiSohnius1975} constrain the mixing
of internal and spacetime symmetries in any relativistic quantum
field theory with a mass gap: the full symmetry group must
factorise as (internal) $\times$ (Poincar\'{e}), up to possible
supersymmetric extensions.  Identifying the internal $3{+}1$
charge-sector grading selected by Theorem~\ref{thm:selection} with
the $3{+}1$ Lorentzian signature of spacetime would constitute
precisely such a mixing, and is therefore not permissible without
additional structure.

\paragraph{What would be needed to bridge the gap.}
Three routes could, in principle, connect the internal grading to
spacetime signature without violating Coleman--Mandula:
\begin{enumerate}\itemsep0pt
  \item \emph{Soldering construction.}  In the Palatini formulation,
    the tetrad $e^a_\mu$ ``solders'' the internal $\sln{4}$ indices
    to spacetime indices.  If the soldering form is present, the
    internal grading induces a spacetime grading.  But this requires
    the soldering form as additional input --- it is not produced by
    the algebraic minimisation alone.
  \item \emph{Pre-geometric regime.}  If the construction operates
    at a level where Poincar\'{e} symmetry has not yet emerged (e.g.,
    a pre-geometric or background-independent formulation), the
    Coleman--Mandula hypotheses do not apply and the internal/spacetime
    distinction is not yet defined.  The grading could then be
    interpreted as a \emph{proto-signature} that seeds the eventual
    spacetime structure upon emergence of a geometric phase.
  \item \emph{Explicit evasion.}  If the theory is formulated with
    massless particles only, or in a way that violates other
    Coleman--Mandula hypotheses (e.g., infinite-dimensional symmetry,
    absence of a mass gap), the no-go theorem does not apply and
    the internal-spacetime identification could be direct.
\end{enumerate}
None of these routes is established in the present work.  The
algebraic result (Theorem~\ref{thm:selection}) stands on its own as
a grading-selection theorem; the physical interpretation as spacetime
signature selection remains an open conjecture requiring one of the
above bridging mechanisms.

\paragraph{The honest scope of the result.}
The theorem selects a $3{+}1$ grading of the internal carrier space
of $\sln{4}$ associated with the $B{-}L$ generator.  This grading
has the same shape as the Lorentzian signature of spacetime.  Whether
this coincidence reflects a deep structural connection (via one of
the bridging mechanisms above) or is merely a numerical coincidence
of the $n = 4$ case is not determined by the algebraic argument
alone.

In coordinate-free language, the mathematical content is:
\begin{quote}
\emph{The minimiser is the grading that makes all active
cross-cluster root spaces odd.  For $\TBL$, this grading has
$(3,1)$ shape.}
\end{quote}

\paragraph{Compact versus noncompact algebras.}
The mechanism extends beyond the $A_n$~series.  Numerical tests
on compact orthogonal algebras $\mathfrak{so}(n)$ ($n = 3, \ldots, 6$),
the noncompact Lorentz algebra $\mathfrak{so}(3,1)$, and the
symplectic algebra $\mathfrak{sp}(4,\mathbb{R})$ confirm the
general pattern (16 cases total, all consistent):
\begin{itemize}\itemsep0pt
  \item \emph{Compact} algebras have $\ad_T^2 < 0$ on the active
    sector.  The minimiser assigns all active generators to the
    even subspace.
  \item \emph{Noncompact} algebras have $\ad_T^2 > 0$ on the active
    sector (or mixed signs in algebras like $\mathfrak{so}(3,1)$).
    The minimiser assigns positive-$\ad_T^2$ generators to the odd
    subspace, producing a nontrivial grading.
\end{itemize}
This parallels the standard relationship between real forms and
the Cartan decomposition $\mathfrak{g} = \mathfrak{k}
\oplus \mathfrak{p}$; whether the two
decompositions coincide in general is an open question.

\paragraph{What is not claimed.}
\begin{enumerate}\itemsep0pt
  \item The result selects a grading of an \emph{internal} carrier
    space, not of physical spacetime.  Bridging the two requires
    additional structure (soldering, pre-geometry, or Coleman--Mandula
    evasion) not provided here.
  \item The result does not explain SM parity violation ($SU(2)_L$
    left-handedness).  That requires the full Pati--Salam product
    group $SU(4)_C \times SU(2)_L \times SU(2)_R$, not
    $\sln{4}$ alone.
  \item The minimisation principle $\delta\widetilde{\mathcal{S}}_g = 0$ is
    proposed, not derived from a deeper action.  It is the simplest
    consistent choice, but ``simplest'' is a criterion, not a proof.
  \item The mechanism has been verified across all four classical
    types ($A_n$, $B_n$, $C_n$, $D_n$) and is dimension-independent.
    Extension to exceptional algebras ($G_2$, $F_4$, $E_{6,7,8}$)
    remains open.
\end{enumerate}

\begin{remark}[Non-commutativity interpretation]
The selection mechanism is equivalent to assigning negative
grading to the minimal index set that maximises the measure of
non-commuting directions under $\ad_T$.
\end{remark}

%-----------------------------------------------------------------
\section*{Acknowledgment}
%-----------------------------------------------------------------

This result emerged from an extended investigation of chirality
structure in the ACS bracket algebra.  An initial measurement of a
``$72\%$ parity-odd holonomy bias'' was classified as
grading-induced after invariance testing.  The upgrade to a
grading-selection theorem occurred when the spectral functional
$\widetilde{\mathcal{S}}_g = \Tr(\sigma_P \cdot g(\ad_T))$ was
found to select a unique $(3,1)$ grading.  Subsequent universality
testing (six functionals, frame invariance, generator dependence,
higher-dimensional generalisation) confirmed the robustness of the
selection.  A physicist's review identified the need to restrict to
$g(0) = 0$ for a clean theorem statement, to separate the $f(0) > 0$
case as a distinct proposition, and --- most importantly --- to
distinguish the internal grading result from any claim about
physical spacetime signature, in light of the Coleman--Mandula
theorem.  These corrections significantly strengthened the paper.

%-----------------------------------------------------------------
\begin{thebibliography}{9}

\bibitem{WallaceA}
B.~Wallace,
\textit{Colour from Gravity: Pati--Salam Structure from the Palatini
Bracket},
2026.

\bibitem{WallaceN1}
B.~Wallace,
\textit{Pythagorean Arithmetic Structure in the Minimal Pati--Salam
Bracket Algebra},
2026.

\bibitem{ColemanMandula1967}
S.~Coleman and J.~Mandula,
\textit{All Possible Symmetries of the S Matrix},
Phys.\ Rev.\ \textbf{159}, 1251 (1967).

\bibitem{HaagLopuszanskiSohnius1975}
R.~Haag, J.~T.~{\L}opusza\'{n}ski, and M.~Sohnius,
\textit{All Possible Generators of Supersymmetries of the S-Matrix},
Nucl.\ Phys.\ B \textbf{88}, 257 (1975).

\end{thebibliography}

\end{document}
